% !TeX program = lualatex
% Copyright (C) 2026 IEEE UFG Computer Society contributors.
%
% This work may be distributed and/or modified under the conditions of the
% LaTeX Project Public License, either version 1.3c of this license or (at your
% option) any later version. The latest version is available at
% https://www.latex-project.org/lppl.txt.
%
% This work has the LPPL maintenance status `maintained'. The Current
% Maintainer is the IEEE UFG Computer Society. The work consists of the files
% listed in MANIFEST.md. See LICENSE for details and artwork exclusions.
%
% GUIA DE EDIÇÃO
% Mantenha exatamente uma opção ativa em cada grupo abaixo. Para trocar uma
% opção, adicione % à linha atual e remova % da alternativa desejada.
\documentclass[
  paper-size=a4,
  % paper-size=letter,
  % font-size=10pt,
  font-size=11pt,
  % font-size=12pt,
  two-sided=true,
  % two-sided=false,
  title-page=true,
  % title-page=false,
  section-numbering=plain,
  % section-numbering=two-digit,
  inner-margin=22mm,
  outer-margin=22mm,
  top-margin=24mm,
  bottom-margin=28mm,
  binding-offset=0mm,
  color-model=rgb
  % color-model=cmyk
]{ieeeufgcs}
\usepackage[brazilian]{babel}
\usepackage{indentfirst}
\usepackage[hidelinks,hypertexnames=false]{hyperref}
% Mantém o espaçamento vertical regular em páginas com tabelas longas.
\raggedbottom

% Para usar referências bibliográficas, renomeie
% references.bib.example para references.bib e ative as três linhas abaixo.
% \usepackage{csquotes}
% \usepackage[backend=biber,style=numeric]{biblatex}
% \addbibresource{references.bib}

% COMANDOS ÚTEIS
% Imagem: \includegraphics[width=.8\linewidth]{img/arquivo.png}
% Referência: \cite{chave-da-referencia}
% Tabela simples:
% \begin{tabular}{ll}
%   Primeira coluna & Segunda coluna \\
% \end{tabular}
% Para tabelas institucionais, copie um dos exemplos IEEECSBrandedTable
% presentes no corpo deste documento.

% Edite somente os valores depois de cada sinal de igual.
\IEEECSSetup{
  document-type={Plano de trabalho},
  title={Título do plano},
  subtitle={Período ou ciclo de referência},
  running-title={Plano de trabalho},
  institution={Capítulo Estudantil da IEEE Computer Society na UFG},
  unit={Unidade acadêmica ou organização responsável},
  author={Equipe responsável},
  date={\today},
  location={Goiânia},
  identifier={Versão 1.0}
}

\begin{document}
\maketitle

\section*{Resumo}
\phantomsection\addcontentsline{toc}{section}{Resumo}

Sintetize a finalidade do plano, o período de execução, as principais
entregas e a decisão institucional solicitada. Substitua este texto antes
da submissão.

\clearpage
\tableofcontents
\clearpage

\section{Apresentação}\label{sec:apresentacao}

Apresente a organização, o contexto do plano e sua relação com a comunidade.
Delimite o período, o público e o escopo institucional do documento.

\section{Motivação e justificativa}

Descreva o problema ou a oportunidade que justifica o plano. Relacione a
proposta com necessidades observáveis e com a missão da organização.
Demonstre que a escala proposta é compatível com os recursos disponíveis.

\section{Membros e público-alvo}

Identifique quem participa da execução e quem será atendido. Informe como
novos membros serão acolhidos, acompanhados e mantidos ativos. Registre
metas de adesão apenas quando houver meios de verificá-las.

\section{Missão e objetivos}

\subsection{Missão}

Declare a finalidade permanente da iniciativa e os princípios que orientam
sua execução.

\subsection{Objetivo geral}

Defina o resultado amplo que o plano pretende alcançar.

\subsection{Objetivos específicos}

\begin{itemize}
  \item converter a missão em resultados observáveis;
  \item estabelecer práticas documentadas e repetíveis;
  \item desenvolver responsáveis e pessoas sucessoras.
\end{itemize}

\section{Governança e equipe responsável}

Descreva a composição da equipe, o processo de decisão, a periodicidade das
reuniões e a relação com as instâncias institucionais. Use a tabela para
registrar a estrutura de responsabilidades e o vínculo entre funções.

\begin{IEEECSBrandedTable}{
  >{\raggedright\arraybackslash}m{0.24\linewidth}
  >{\raggedright\arraybackslash}X
  >{\raggedright\arraybackslash}m{0.25\linewidth}
}
\IEEECSHeaderCell{Função} &
\IEEECSHeaderCell{Responsabilidade} &
\IEEECSHeaderCell{Responsável} \\
Coordenação & Aprovar prioridades e acompanhar a execução. & A definir. \\
Secretaria & Registrar decisões e manter documentos. & A definir. \\
Responsabilidade técnica & Orientar as entregas e revisar riscos. &
A definir. \\
Orientação docente & Acompanhar o plano e a equipe. & A definir. \\
\end{IEEECSBrandedTable}

\subsection{Perfis dos integrantes}

Registre os perfis das pessoas que assumem funções no plano. Duplique a
tabela para cada integrante e substitua os campos pelo conteúdo pertinente.
Evite informações pessoais que não sejam necessárias à avaliação do plano.

\begin{IEEECSBrandedTable}{
  >{\raggedright\arraybackslash}m{0.28\linewidth}
  >{\raggedright\arraybackslash}X
}
\IEEECSHeaderCell{Campo} & \IEEECSHeaderCell{Informação} \\
Nome & Nome completo. \\
Função & Cargo ou responsabilidade no plano. \\
Vínculo & Unidade acadêmica ou organização. \\
Experiência relevante & Competências ligadas às atividades previstas. \\
Contato institucional & \url{https://example.org/contato} \\
\end{IEEECSBrandedTable}

\subsection{Sucessão}

Explique como responsabilidades, documentos, acessos e conhecimento serão
transferidos. Identifique responsáveis principais e pessoas em formação.

\section{Projetos e atividades}

Organize as frentes de trabalho e associe cada uma a uma entrega mínima.
Use a Seção~\ref{sec:cronograma} para indicar quando as atividades ocorrerão.

\begin{IEEECSBrandedTable}{
  >{\raggedright\arraybackslash}m{0.24\linewidth}
  >{\raggedright\arraybackslash}X
  >{\raggedright\arraybackslash}X
}
\IEEECSHeaderCell{Frente} &
\IEEECSHeaderCell{Propósito} &
\IEEECSHeaderCell{Primeiro ciclo} \\
Formação & Desenvolver competências técnicas e organizacionais. &
Atividades e evidências de aprendizagem. \\
Projetos & Transformar necessidades delimitadas em entregas verificáveis. &
Protótipo ou relatório documentado. \\
Comunidade & Integrar participantes e manter comunicação recorrente. &
Calendário, acolhimento e evento de referência. \\
\end{IEEECSBrandedTable}

\subsection{Escalabilidade}

Defina critérios de passagem antes de ampliar o público, os recursos ou a
complexidade operacional.

\begin{IEEECSBrandedTable}{
  >{\raggedright\arraybackslash}m{0.20\linewidth}
  >{\raggedright\arraybackslash}X
  >{\raggedright\arraybackslash}X
}
\IEEECSHeaderCell{Etapa} &
\IEEECSHeaderCell{Critério de passagem} &
\IEEECSHeaderCell{Entrega principal} \\
Fundação & Responsáveis, regras e recursos iniciais confirmados. &
Plano aprovado e calendário inicial. \\
Piloto & Primeiras atividades testadas e carga da equipe conhecida. &
Piloto documentado e retrospectiva. \\
Consolidação & Rotina estável e sucessão preparada. &
Relatório do ciclo e transição de gestão. \\
\end{IEEECSBrandedTable}

\section{Cronograma}\label{sec:cronograma}

Substitua os períodos de exemplo por datas compatíveis com o calendário
institucional. Indique os responsáveis e as evidências de conclusão.

\begin{IEEECSBrandedTable}{
  >{\raggedright\arraybackslash}m{0.17\linewidth}
  >{\raggedright\arraybackslash}m{0.23\linewidth}
  >{\raggedright\arraybackslash}X
  >{\raggedright\arraybackslash}m{0.20\linewidth}
}
\IEEECSHeaderCell{Período} &
\IEEECSHeaderCell{Prioridade} &
\IEEECSHeaderCell{Entrega e evidência} &
\IEEECSHeaderCell{Responsável} \\
M1--M3 & Fundar e testar & Plano aprovado e primeiro piloto registrado. &
Coordenação. \\
M4--M6 & Validar e abrir & Materiais revisados e participação ampliada. &
Equipe de formação. \\
M7--M9 & Integrar & Projetos executados e resultados documentados. &
Lideranças de projeto. \\
M10--M12 & Transferir & Relatório do ciclo e sucessão preparados. &
Coordenação e secretaria. \\
\end{IEEECSBrandedTable}

\section{Metas e indicadores}\label{sec:metas}

Defina metas verificáveis para o período. Registre como cada resultado será
medido e quem acompanhará a informação.

\begin{IEEECSNotice}[Condição de validade]
Nenhuma meta deve prevalecer sobre segurança, integridade, acessibilidade,
capacidade das pessoas voluntárias ou disponibilidade real de recursos.
\end{IEEECSNotice}

\begin{IEEECSBrandedTable}{
  >{\raggedright\arraybackslash}m{0.22\linewidth}
  >{\raggedright\arraybackslash}X
  >{\raggedright\arraybackslash}m{0.18\linewidth}
  >{\raggedright\arraybackslash}m{0.17\linewidth}
}
\IEEECSHeaderCell{Meta} &
\IEEECSHeaderCell{Indicador e fonte} &
\IEEECSHeaderCell{Prazo} &
\IEEECSHeaderCell{Responsável} \\
Realizar o piloto & Relatório de execução aprovado. & M3 & Coordenação. \\
Formar participantes & Lista de presença e avaliação da atividade. & M6 &
Equipe de formação. \\
Documentar o ciclo & Relatório anual entregue à instância responsável. &
M12 & Secretaria. \\
\end{IEEECSBrandedTable}

\section{Parcerias e relações institucionais}

Identifique as unidades acadêmicas, organizações e comunidades que poderão
contribuir. Para cada parceria, indique o objetivo, o contato responsável e
a condição necessária para formalizá-la. Diferencie contatos em negociação
de compromissos já confirmados.

\section{Recursos e sustentabilidade}

Estime os recursos humanos, físicos, digitais e financeiros necessários.
Indique fontes confirmadas, fontes potenciais e como eventuais despesas
serão aprovadas e registradas.

\begin{IEEECSBrandedTable}{
  >{\raggedright\arraybackslash}m{0.28\linewidth}
  >{\raggedright\arraybackslash}X
  >{\raggedright\arraybackslash}m{0.24\linewidth}
}
\IEEECSHeaderCell{Recurso} &
\IEEECSHeaderCell{Necessidade e finalidade} &
\IEEECSHeaderCell{Fonte ou situação} \\
Espaço físico & Local para reuniões e atividades. & A confirmar. \\
Plataforma digital & Registro e compartilhamento de materiais. &
Disponível ou a confirmar. \\
Orçamento & Custos previstos por atividade. & Estimar e aprovar. \\
\end{IEEECSBrandedTable}

\section{Comunicação e divulgação}

Indique os canais institucionais, o público de cada canal, a frequência de
atualização e quem revisará as publicações. Registre como serão divulgados
convites, resultados e mudanças no cronograma.

\section{Monitoramento e relatórios}

Acompanhe periodicamente as metas da Seção~\ref{sec:metas}. Registre
as reuniões, guarde as evidências e defina o formato dos relatórios.
Indique quem consolidará os resultados e a quem serão apresentados.

\section{Riscos e mitigação}

Identifique os riscos mais relevantes para este plano. Revise probabilidade,
impacto e medidas de resposta ao longo da execução.

\begin{IEEECSBrandedTable}{
  >{\raggedright\arraybackslash}m{0.21\linewidth}
  >{\raggedright\arraybackslash}m{0.14\linewidth}
  >{\raggedright\arraybackslash}m{0.13\linewidth}
  >{\raggedright\arraybackslash}X
}
\IEEECSHeaderCell{Risco} &
\IEEECSHeaderCell{Probabilidade} &
\IEEECSHeaderCell{Impacto} &
\IEEECSHeaderCell{Mitigação} \\
Sobrecarga da equipe & Média & Alto &
Revisar a carga de trabalho e reduzir o escopo quando necessário. \\
Atraso em recursos & Média & Médio &
Confirmar dependências antes de agendar cada atividade. \\
Baixa adesão & Média & Médio &
Revisar a divulgação e ouvir o público-alvo. \\
\end{IEEECSBrandedTable}

\section{Pedido de aprovação}

Registre de maneira direta quais decisões são solicitadas.

\begin{enumerate}
  \item aprovar o plano e seu período inicial;
  \item confirmar responsáveis e condições institucionais;
  \item autorizar as primeiras atividades previstas.
\end{enumerate}

\subsection{Primeiros passos após a aprovação}

Liste as ações imediatas, seus responsáveis e os prazos iniciais.

\section{Conclusão}

Retome o valor esperado, as condições de execução e a decisão solicitada.

\clearpage
\appendix

\section{Documentos e registros de apoio}

Liste minutas, formulários, atas ou registros que acompanham o plano. Indique
onde cada item será mantido. Remova este apêndice quando não for necessário.

\section{Termo de compromisso e assinaturas}

Descreva o compromisso assumido pelas pessoas responsáveis. Insira as
assinaturas eletrônicas no documento final depois da aprovação do texto.

\begin{IEEECSBrandedTable}{
  >{\raggedright\arraybackslash}m{0.38\linewidth}
  >{\raggedright\arraybackslash}X
}
\IEEECSHeaderCell{Nome e função} &
\IEEECSHeaderCell{Assinatura eletrônica e data} \\
Nome completo --- função & Espaço reservado para assinatura e data. \\
Nome completo --- orientação docente &
Espaço reservado para assinatura e data. \\
\end{IEEECSBrandedTable}

% Ative esta linha depois de configurar o arquivo references.bib.
% \printbibliography

\end{document}
